\section{The five pair-sum patterns}
\label{sec:classification}

The ordinary-convolution constraints depend only on equalities between
unordered pair sums, with repeated summands allowed.  For four integer
points these equalities give five patterns.

\begin{proposition}\label{c-classification}
Every four-point integer set belongs to exactly one of the following
classes: a Sidon set, an isolated three-term arithmetic progression, a
proper parallelogram, two three-term arithmetic progressions, or a
four-term arithmetic progression.  Here a Sidon set has distinct unordered
pair sums.  An isolated progression has exactly one nontrivial pair-sum
equality, of the form $p+r=2q$.  A proper parallelogram has exactly one
nontrivial pair-sum equality, with four distinct summands.

Up to translation, reflection, and integer dilation, the class with two
progressions is represented by $\{0,1,2,4\}$.
\end{proposition}

\begin{proof}
Translate the set to $\{0,a,b,c\}$, where $0<a<b<c$.  A shared summand
can be cancelled from any pair-sum equality.  The remaining nontrivial
possibilities are precisely
\begin{equation}\label{c-relations}
 b=2a,\qquad c=2a,\qquad c=2b,\qquad
 c=2b-a,\qquad c=a+b.
\end{equation}
The first four assert that one point is the midpoint of two others.
The last is the only possible equality between pairs using all four points.

The compatible intersections are as follows.  The equations $b=2a$ and
$c=2b$ give $\{0,a,2a,4a\}$.  The equations $c=2a$ and $c=2b-a$ give
$b=3a/2$, hence a reflected copy of $\{0,d,2d,4d\}$ with $d=a/2$.
Any two of $b=2a$, $c=2b-a$, and $c=a+b$ imply the third and give
$\{0,a,2a,3a\}$.  Every other pair of equations in
\eqref{c-relations} contradicts $0<a<b<c$.  The possibilities are
therefore no equality, one midpoint equality, one cross-pair equality,
the two-progressions case, and the four-term progression.
\end{proof}

This classification concerns pair sums.  A class can contain further
equalities between fourfold sums, which must also be accounted for when
bounding $T_F(k)$.

\section{Sidon supports}
\label{sec:sidon}

\begin{lemma}\label{c-level-set}
For every finitely supported nonnegative function $h$ on $\mathbb Z$,
\begin{equation}\label{c-level-bound}
 \|h\star h\|_1\ge 2\sum_s h(s)^2-\|h\|_\infty^2.
\end{equation}
\end{lemma}

\begin{proof}
The zero function is immediate.  For $0\le t<\|h\|_\infty^2$, put
$A_t=\{s:h(s)^2>t\}$.  Then $A_t$ is nonempty and
$A_t+A_t\subseteq\{s:(h\star h)(s)>t\}$.  If a finite integer set has
ordered elements $a_1<\cdots<a_m$, the chain
\[
 2a_1<a_1+a_2<\cdots<a_1+a_m<a_2+a_m<\cdots<2a_m
\]
shows that $|A+A|\ge2|A|-1$.  Integrating the resulting inequality for
superlevel-set cardinalities over $0\le t<\|h\|_\infty^2$ proves
\eqref{c-level-bound}.
\end{proof}

\begin{proposition}\label{c-sidon}
Let $F$ be a finite nonempty Sidon set and let $k\ge0$ on $F$.  Put
$K=\sum_{p\in F}k(p)$ and $a=\max_{p\in F}k(p)$.  Then
\begin{equation}\label{c-sidon-norm}
 N_F(k)=\max\left\{\frac K2,\frac{K+3a}{4}\right\}.
\end{equation}
If $|F|\le4$, then $N_F(k)^4\le T_F(k)$, with strict inequality
whenever at least two values of $k$ are positive.
\end{proposition}

\begin{proof}
The Sidon property makes the ordinary-convolution constraints equivalent to
\[
 w_p\le1,\qquad w_pw_q\le\frac14\quad(p\ne q),
 \qquad w_p=x(p)^2\ge0.
\]
If all $w_p\le1/2$, the objective is at most $K/2$.  Otherwise let
$w_j=u\in[1/2,1]$ be largest.  Every other coordinate is at most
$1/(4u)$, so
\[
 \sum_p k(p)w_p\le k(j)u+\frac{K-k(j)}{4u}.
\]
This is convex in $u$, and its endpoint values are $K/2$ and
$k(j)+(K-k(j))/4$.  Maximizing over $j$ proves the upper bound in
\eqref{c-sidon-norm}.  The profiles with all coordinates $1/2$, and
with one coordinate $1$ and all others $1/4$, are feasible.  They attain
the stated values.

Suppose now that $|F|\le4$.  The case $k=0$ is immediate.  Apply
Lemma~\ref{c-level-set} to $h=k\star k$.  Distinct unordered pair sums
give
\[
 \sum_s h(s)^2
 =\frac12\left[\left(\sum_p k(p)^2\right)^2+\sum_p k(p)^4\right],
 \qquad \|h\|_\infty=a^2.
\]
Associativity of max-convolution and Cauchy--Schwarz therefore imply
\begin{equation}\label{c-sidon-half}
 T_F(k)\ge
 \left(\sum_p k(p)^2\right)^2+\sum_p k(p)^4-a^4
 \ge\left(\frac K2\right)^4.
\end{equation}
The last comparison is strict if at least two kernel values are positive,
because then $\sum_p k(p)^4-a^4>0$.

For the other term in \eqref{c-sidon-norm}, choose a point $p_0$ with
$k(p_0)=a$ and write $R=K-a$.  Fixing two copies of $p_0$ in the
fourfold maximum gives distinct sums indexed by unordered pairs from $F$.
Thus
\begin{equation}\label{c-pivot-bound}
 T_F(k)\ge a^4+a^3R+a^2h_2
 \ge a^4+a^3R+\frac23a^2R^2,
\end{equation}
where $h_2$ is the sum of the squares and pair products of the other
kernel values.  There are at most three such values, so
$h_2=(R^2+\sum_{p\ne p_0}k(p)^2)/2\ge2R^2/3$.
Since $R\le3a$, expansion gives
\begin{equation}\label{c-quarter-bound}
 \left(a+\frac R4\right)^4
 \le a^4+a^3R+\frac{153}{256}a^2R^2.
\end{equation}
Indeed the coefficient after bounding the cubic and quartic terms is
$3/8+3/16+9/256=153/256$.  The difference between the right sides of
\eqref{c-pivot-bound} and \eqref{c-quarter-bound} is at least
$53a^2R^2/768$, positive when $R>0$.  This proves both comparisons in
\eqref{c-sidon-norm}, including the asserted strictness.
\end{proof}

\section{Proper parallelograms}
\label{sec:parallelogram}

Write
\begin{equation}\label{c-parallelogram-set}
 F=\{0,a,b,a+b\},\qquad 0<a<b,\qquad b\ne2a,
\end{equation}
and index its coordinates in this order.  The opposite pairs are
$\{0,3\}$ and $\{1,2\}$.  Define a collection $\mathcal P$ of sixteen
squared profiles:
\begin{itemize}
 \item Four profiles have three coordinates $1/2$ and one coordinate zero.
 \item Eight profiles have one coordinate $1$, its opposite coordinate
       and one adjacent coordinate $1/4$, and the remaining coordinate zero.
 \item Four profiles have one coordinate $1$, its two adjacent coordinates
       $1/4$, and its opposite coordinate $1/16$.
\end{itemize}
We call the last four profiles tensor profiles, since their entries are
products of the two-coordinate vector $(1,1/4)$.

\begin{proposition}\label{c-parallelogram}
For every $k\ge0$ on \eqref{c-parallelogram-set},
\begin{equation}\label{c-parallelogram-norm}
 N_F(k)=\max_{v\in\mathcal P}k\cdot v,
 \qquad (k\cdot v)^4\le T_F(k)\quad(v\in\mathcal P).
\end{equation}
If $k$ is positive at all four points, every displayed quartic comparison
is strict, and hence $N_F(k)^4<T_F(k)$.
\end{proposition}

\subsection{Reduction to sixteen profiles}

Ordinary-convolution feasibility is equivalent to
\begin{equation}\label{c-parallelogram-feasible}
 x_i^2\le1,\qquad 2x_ix_j\le1\ (i\ne j),\qquad
 x_0x_3+x_1x_2\le\frac12.
\end{equation}
The last constraint is the coefficient at $a+b$; its separate opposite-pair
bounds are redundant.  Each profile in $\mathcal P$ is feasible.
For the first two types the sum of the two opposite products of square
roots is $1/2$; for a tensor profile both products are $1/4$.

Put $w_i=x_i^2$.  First suppose every $w_i\le1/2$.  With
$r=x_0x_3$ and $s=x_1x_2$, the inequality
$u+v\le1/2+2uv$ for $u,v\in[0,1/2]$ gives
\[
 \sum_iw_i\le1+2r^2+2s^2\le1+2(r+s)^2\le\frac32.
\]
Increase coordinates within $[0,1/2]$ until their sum is $3/2$.
The resulting vector lies in the convex hull of the four profiles of the
first type: for such a vector $y$, the coefficient of the profile with its
$i$th coordinate zero is $1-2y_i$.  These coefficients are nonnegative
and sum to one.

Otherwise relabel corners, preserving opposite pairs, so that $w_0=u\ge1/2$,
and put $L=1/(4u)$.  Then $w_1,w_2,w_3\le L$, and
\[
 w_3\le L\left(1-2\sqrt{w_1w_2}\right)^2.
\]
For fixed $u$, the upper bound
\[
 k_0u+k_1w_1+k_2w_2
       +k_3L\left(1-2\sqrt{w_1w_2}\right)^2
\]
is separately convex in $w_1,w_2\in[0,L]$.  Its only nonlinear term in
either variable is a nonpositive multiple of that variable's square root.
It is therefore enough to consider the four corners of this square.

The corner $(0,0)$ is dominated by either $(L,0)$ or $(0,L)$.  These
give $(u,L,0,L)$ and $(u,0,L,L)$, whose objectives are convex in $u$.
Their endpoint values correspond to profiles of the first two types.
The last corner gives
\[
 b(u)=\left(u,L,L,\frac{(2u-1)^2}{16u^3}\right).
\]
For this curve, set
\[
 t=2u-1,\qquad \nu=1-t,\qquad
 s=\frac{(2u-1)(1-u)}{4u},
\]
and define
\[
 U=(1,1/4,1/4,1/16),\quad W_0=(1/2,1/2,1/2,0),
\]
\[
 W_1=(1/2,1/2,0,1/2),\qquad W_2=(1/2,0,1/2,1/2).
\]
The coefficients of
\[
 y=tU+(\nu-4s)W_0+2sW_1+2sW_2
\]
are nonnegative and sum to one, since $4s=t(1-u)/u\le2(1-u)=\nu$.
The first three coordinates of $y$ are $(u,L,L)$ and its last coordinate
is $t/16+2s$.  Moreover,
\[
 \frac{(2u-1)^2}{16u^3}-\frac t{16}
 =\frac{t(1-u)(u^2+u-1)}{16u^3}
 \le\frac s4\le2s.
\]
Thus $y$ dominates $b(u)$.  These reductions prove the upper bound in
\eqref{c-parallelogram-norm}; feasibility of the listed profiles proves
equality.

\subsection{Profiles supported on three corners}

Every three-corner subset of a proper parallelogram is Sidon.  An additional
midpoint equality would force $b=2a$ or $a=b$.  For a kernel restricted to
such a triangle, Proposition~\ref{c-sidon} bounds both its half-sum and every
unit-and-two-quarter form.  These are exactly the objectives of the twelve
profiles with a zero coordinate.  Monotonicity under enlarging the support
then proves their quartic comparisons with $T_F(k)$.  For positive $k$,
their three-point restrictions have at least two positive values, so the
comparisons are strict.

\subsection{Tensor profiles and their collision grid}

Reflection around $(a+b)/2$ reduces the four tensor profiles to a pivot at
$0$ and a pivot at $a$.  Relative to the pivot, write the kernel values
as $A,B,C,D$, with the tensor objective
\[
 E=A+B/4+C/4+D/16.
\]
For the outer pivot the offsets are $0,a,b,a+b$.  For the inner pivot
they are $0,-a,b,b-a$, so that $A=k(a)$, $B=k(0)$, $C=k(a+b)$, and
$D=k(b)$ in that case.

A monomial with multiplicities $n_A,n_B,n_C,n_D$ has grid coordinates
$i=n_B+n_D$ and $j=n_C+n_D$.  Both belong to $\{0,1,2,3,4\}$.
The total coefficient of its grid point in $E^4$ is
\begin{equation}\label{c-grid-weights}
 c_ic_j,\qquad (c_0,c_1,c_2,c_3,c_4)=(1,1,3/8,1/16,1/256),
\end{equation}
because the coefficient-generating polynomial is
\[
 (1+z/4+w/4+zw/16)^4=(1+z/4)^4(1+w/4)^4.
\]
The actual sum is $ia+jb$ for the outer pivot and $-ia+jb$ for the inner
pivot.  Grouping by that sum proves $E^4\le T_F(k)$ whenever every
group's nonnegative coefficient weights add to at most one.  AM--GM
inequalities can first move terms between groups, provided each replacement
is a legal fourfold monomial at its new sum.

Write $b/a=r/s$ in lowest terms, with $r>s>0$.  Two grid points in
$\{0,\ldots,4\}^2$ collide only if $r\le4$.  Excluding $r/s=2$ leaves
\begin{equation}\label{c-exceptional-ratios}
 r/s\in\{3,4,3/2,4/3\}.
\end{equation}
The collision step is $(r,-s)$ for the outer pivot and $(r,s)$ for the
inner pivot.  Since $r\ge3$, no group contains three grid points.
Outside \eqref{c-exceptional-ratios}, \eqref{c-grid-weights} proves the
quartic comparison directly.

We use two elementary inequalities:
\begin{align}
 U+\epsilon V&\le\max(U,V)+\epsilon\min(U,V)
     &&(U,V\ge0,\ 0\le\epsilon\le1),\label{c-max-min}\\
 \min(A^3C,X^4)&\le A^2X^2+A^2C^2
     &&(A,C,X\ge0).\label{c-absorption}
\end{align}
The first follows by ordering $U,V$.  For the second, the case $A=0$
is immediate.  Otherwise divide by $A^4$ and put $x=X/A$, $y=C/A$.
If $x\le1$, use $\min(y,x^4)\le x^4\le x^2$; if $x\ge1$ and
$y\le1$, use $y\le x^2$; if $y\ge1$, use $y\le y^2$.

\subsection{Outer-pivot collisions}

At ratios $3/2$ and $4/3$, every colliding pair has total weight at most
$3/8+1/16=7/16$.  Indeed the point with smaller first coordinate has
second coordinate at least two, and the other point has first coordinate
at least three.  Thus the coefficient comparison applies directly.

At ratio three, the only groups of weight greater than one are
$(0,1)\sim(3,0)$ and $(1,1)\sim(4,0)$.  Their polynomials are
\[
 A^3C+AB^3/16,\qquad
 3A^2BC/4+A^3D/4+B^4/256.
\]
Use $AB^3/16\le(A^2B^2+B^4)/32$ in the first group.  The resulting
total coefficient of $B^4$ in the second group is $\delta=9/256$.
Keep
\[
 (3/4-\delta)A^2BC+A^3D/4+\delta B^4
\]
there; its coefficient sum is one.  The displaced term satisfies
$\delta A^2BC\le\delta(A^2B^2+A^2C^2)/2$.
The extra coefficients assigned to $A^2B^2$ and $A^2C^2$ are
$25/512$ and $9/512$.  Their entire original group weights are
$3/8$ and $7/16$, respectively, so the final weights are
$217/512$ and $233/512$.  Both are below one.

At ratio four, the sole overloaded group is $A^3C+B^4/256$.
Equations~\eqref{c-max-min}--\eqref{c-absorption}, with
$\epsilon=1/256$ and $X=B$, bound it by its maximum plus
$(A^2B^2+A^2C^2)/256$.  The two recipient groups originally have weights
$3/8$ and $97/256$; their final weights are $97/256$ and $98/256$.
Every other group already has weight at most one.

\subsection{Inner-pivot collisions}

An inner-pivot group can be overloaded only if the smaller grid point has
both coordinates at most one.  Otherwise its weight is at most $3/8$,
and that of its partner is at most $1/16$.
For ratios three, four, and $3/2$, remove the high-point polynomials in
the following table and keep their low-point partners in full.
\begin{center}
\begin{tabular}{cl}
\toprule
High grid point & Polynomial \\
\midrule
$(3,1)$ & $B^3C/64+3AB^2D/64$ \\
$(4,1)$ & $B^3D/256$ \\
$(3,2)$ & $3B^2CD/256+3ABD^2/256$ \\
$(4,2)$ & $3B^2D^2/2048$ \\
$(3,3)$ & $3BCD^2/1024+AD^3/1024$ \\
$(4,3)$ & $BD^3/4096$ \\
\bottomrule
\end{tabular}
\end{center}
The entries follow from the multinomial expansion of $E^4$.  The sets
of removed points and their total coefficient weights are
\begin{center}
\begin{tabular}{clc}
\toprule
Ratio & Removed points & Total weight \\
\midrule
$3$ & $(3,1),(4,1),(3,2),(4,2)$ & $187/2048$ \\
$4$ & $(4,1),(4,2)$ & $11/2048$ \\
$3/2$ & $(3,2),(4,2),(3,3),(4,3)$ & $119/4096$ \\
\bottomrule
\end{tabular}
\end{center}
Each total is below $1/8$.  Every unweighted monomial in the first table
is at most
\[
 S=B^4+C^4+D^4+A^2B^2+A^2D^2.
\]
For monomials without $A$, this follows from weighted AM--GM.  The three
remaining forms satisfy
\[
 AB^2D\le\frac{A^2D^2+B^4}{2},\qquad
 ABD^2\le\frac{A^2B^2+D^4}{2},\qquad
 AD^3\le\frac{A^2D^2+D^4}{2}.
\]
The removed polynomial is consequently bounded by $S/8$.

The five recipient monomials occupy distinct groups, outside the removed
groups.  In the order $B^4,C^4,D^4,A^2B^2,A^2D^2$, their original
weights are
\[
 \frac1{256},\quad\frac1{256},\quad
 \begin{cases}
  1/16+1/65536,&r/s=3\text{ or }4,\\
  3/8+1/65536,&r/s=3/2,
 \end{cases}
 \quad\frac38,\quad\frac9{64}.
\]
These are full group weights, including any collision partner.  For example,
the partner of the $D^4$ point $(4,4)$ is $(1,3)$, $(0,3)$, or $(1,2)$
at the respective ratios.  The other recipient grid points are
$(4,0),(0,4),(2,0),(2,2)$ and are unpaired.  Adding $1/8$ to each
weight leaves every total below one, proving all three cases.

At ratio $4/3$, the two overloaded groups are
\[
 A^4+BD^3/4096,\qquad A^3C+D^4/65536.
\]
Weighted AM--GM gives
\[
 BD^3/4096+D^4/65536
 \le B^4/16384+13D^4/65536.
\]
Keep $A^4$ in its group and combine the $D^4$ term with $A^3C$ using
\eqref{c-max-min}--\eqref{c-absorption}, with
$\epsilon=13/65536$ and $X=D$.  This adds $1/16384$ at $B^4$ and
$13/65536$ at each of $A^2D^2,A^2C^2$.  Their distinct recipient
groups have original weights $1/256,9/64,3/8$, so all additions fit.

This proves every tensor comparison.  For positive $A,B,C,D$, it is
strict: the group containing $A^2B^2$ has a positive maximum and final
weight below one.  Its original weight is $3/8$; the outer-pivot transfers
leave it at $217/512$ or $97/256$, the first three inner-pivot transfers
leave it at $1/2$, and the inner ratio $4/3$ leaves it unchanged.
Together with the twelve three-corner profiles, this proves
Proposition~\ref{c-parallelogram}.  The finite maximum in
\eqref{c-parallelogram-norm} preserves strictness.
